\section{Contribution : contraintes cognitives et alignement}

\begin{frame}{Laboratoire}
\begin{itemize}
    \item \textbf{LSCP} : Laboratoire de Sciences Cognitives et Psycholinguistique
    \item ENS-PSL \& CNRS, Paris
\end{itemize}
\vfill
\centering
\includegraphics[height=1cm]{images/lscp.png}\hspace{1cm}\includegraphics[height=1.3cm]{images/ens_psl.png}\hspace{1cm}\includegraphics[height=1cm]{images/cnrs.png}
\end{frame}

\begin{frame}{Question de recherche}
\begin{block}{Les contraintes cognitives rapprochent-elles les modèles de l'humain ?}
    Deux contraintes évaluées : une \textbf{mémoire limitée} et un entraînement sur \textbf{peu de données} (10 à 100 millions de mots).
\end{block}
\begin{itemize}
    \item 118 modèles de langage comparés à des temps de lecture et à des enregistrements cérébraux
    \item Cerveau : fMRI, EEG, MEG et intracrânien
    \item Méthode : analyse de similarité représentationnelle (RSA)
\end{itemize}
\end{frame}

\begin{frame}{Résultats}
\begin{itemize}
    \item Les contraintes augmentent la similarité avec les \textbf{temps de lecture} et la diminuent avec les \textbf{réponses cérébrales}
    \item Les modèles de grande taille présentent le profil inverse
    \item Aucun modèle n'est le meilleur sur les deux cibles
\end{itemize}
\vspace{0.4em}
\begin{block}{Conclusion}
    La similarité comportementale et la similarité cérébrale sont dissociables.
\end{block}
% IMAGE : figure 2 de l'article (comportement vs cerveau), à ajouter
\src{Travail en cours}
\end{frame}
